This section compares OT and CRDTs based on the criteria defined in the methodology.

\subsection{Response Time}

Both OT and CRDTs allow a node to apply changes locally without waiting for other nodes, which means
neither approach introduces noticeable latency for the user. However, OT requires the server to
transform and rebroadcast operations before other clients can apply them, introducing a round-trip
dependency. CRDTs have no such requirement -- each node can merge incoming updates independently
at any point, making them more suitable for fully asynchronous scenarios.

\subsection{Consistency}

Both approaches guarantee eventual consistency, meaning all nodes will converge to the same state
once all updates have been delivered. OT achieves this through its transformation function, provided
the TP1 property is satisfied \cite{nichols1995jupiter}. CRDTs achieve it structurally, by design:
the merge operation is commutative, associative, and idempotent, so the result is always the same
regardless of the order or number of times updates are applied \cite{shapiro2011comprehensive}.
CRDTs therefore offer a stronger correctness guarantee, since their convergence does not depend on
a correctly implemented transformation function.

\subsection{Difficulty of Implementation}

OT is notoriously difficult to implement correctly. The transformation function must handle all
combinations of concurrent operations, and subtle mistakes can cause divergence. Several published
OT algorithms were later found to contain bugs \cite{nichols1995jupiter}. CRDTs are easier to
reason about at the data structure level, since the merge logic is self-contained and mathematically
provable. However, choosing or designing the right CRDT for a given use case still requires a solid
understanding of the underlying properties.

\subsection{Suitability for Decentralised Use}

OT in its most common form relies on a central server to serialize operations and apply
transformations in a consistent order. Without a server, clients would need to implement TP2, which
is significantly more complex and remains an open research problem \cite{ellis1989concurrency}.
CRDTs, by contrast, are naturally suited to peer-to-peer architectures. Since every merge is
deterministic and order-independent, nodes can synchronize directly with each other without any
central coordinator, and can continue working while offline \cite{automerge2024, nicolaescu2015}.

\subsection{Storage and Network Cost}

State-based CRDTs require transmitting the full state on each synchronization, which becomes costly
as the state grows. Delta-based CRDTs reduce this overhead by sending only incremental changes, but
current algorithms can still produce redundant deltas \cite{Stateful_delta}. Operation-based CRDTs
send only individual operations, but require reliable causal delivery \cite{sending_order}.
OT requires the server to retain the full operation history in order to transform late-arriving
operations from clients that are far behind the current version, which imposes a storage overhead
on the server side.

\subsection{Summary}

Table~\ref{tab:comparison} summarizes the comparison across all criteria.

\begin{table}[h]
\caption{Comparison of OT and CRDTs}
\label{tab:comparison}
\begin{tabular}{|l|c|c|}
\hline
\textbf{Criterion} & \textbf{OT} & \textbf{CRDT} \\
\hline
Response time       & Local-first & Local-first \\
\hline
Consistency         & Eventual (TP1) & Strong eventual \\
\hline
Implementation      & Complex & Moderate \\
\hline
Decentralised use   & Limited & Native support \\
\hline
Network cost        & Operation history & State or deltas \\
\hline
\end{tabular}
\end{table}
